\section{Discussion and Limitations}
\label{sec:discussion}

We state honestly what $\Mtype\langle T \rangle$ does not do, and where
adjacent literature handles similar problems differently. The aim is to
prevent claims the type cannot support, and to mark the boundary of the
contribution so that future work has a clear target.

\subsection{The independence assumption}
\label{sec:disc:independence}

The two components of a $\Mtype\langle T \rangle$ value are modelled as
independent random variables. In Rubin's terms this is missing at random
(MAR)~\cite{rubin1976missing}: the event of presence carries no
information about the value the variable would take if observed. The
assumption is a deliberate scope choice. It captures intermittent sensors
and many clinical-trial reporting patterns, and it lets the
implementation reuse the existing $\Utype\langle T \rangle$ sampling and
caching machinery without change.

The assumption fails for missing-not-at-random (MNAR) data, where the
unobserved value influences the probability of being observed. A glucose
sensor that fails preferentially during high readings, a survey
respondent who declines to answer when the answer is socially
undesirable, and a patient who stops reporting because the drug stopped
working all violate independence. For these workloads, the correct
construction is a joint distribution over presence and value, and the
type would need a new constructor accepting that joint distribution
together with new sampling logic that respects the conditional structure.
We mark this as the principal direction for future work.

\subsection{SPRT termination}
\label{sec:disc:sprt}

The gating operator runs an SPRT against a caller-supplied sample cap.
Wald's optimality result~\cite{wald1948optimum} bounds the
\emph{expected} number of samples; it does not bound the worst case.
When the operand's true presence probability sits inside the
indifference region $[p_\tau - \varepsilon, p_\tau + \varepsilon]$, the
test typically runs to the cap before falling back to the empirical
comparison $\bar p > p_\tau$. The benchmark in~\Cref{sec:eval:sprt}
quantifies this regime at a few microseconds with the default cap of
$N = 1000$.

The fallback decision is sound under the empirical evidence collected but
does not carry the same Type I and Type II error guarantees as the
SPRT-accepted outcomes. Callers needing tight worst-case error budgets
should choose $\varepsilon$ and $N$ together, with $\varepsilon$ large
enough to keep the operand out of the indifference region under expected
operating conditions and $N$ large enough that the fallback rate is
acceptable. Group-sequential alternatives~\cite{jennison2000groupseq}
exist that provide closed worst-case bounds; we did not adopt one for
this work but note them as a candidate replacement for the SPRT.

\subsection{Nested presence}
\label{sec:disc:nesting}

Constructions of the form $\Mtype\langle \Mtype\langle T \rangle \rangle$
are not the intended use of the type. The trait bound
\texttt{ProbabilisticType} requires \texttt{Clone}, \texttt{Send},
\texttt{Sync}, and \texttt{IntoSampledValue} / \texttt{FromSampledValue}
on the carrier type; $\Mtype\langle T \rangle$ does not implement these
trait conversions, so the Rust type system refuses the nesting at
compile time. The design rationale is that two layers of probabilistic
presence are almost never what an application means: nested presence
collapses semantically (``the optional optional value'' is just the
optional value), and the nested encoding obscures the single Bernoulli
that the modeller actually cares about. We treat the absence of nesting
as a feature rather than a limitation.

\subsection{Quality of the presence prior}
\label{sec:disc:prior}

The Bernoulli parameter $p$ supplied to $\textsc{fromBernoulli}$ is the
modeller's responsibility. A wrong $p$ poisons the analysis: too high
and the gating operator passes data that should have been excluded; too
low and the operator excludes data that should have been included. The
type does not estimate $p$ for the caller; it accepts what is given. For
intermittent sensors $p$ can be estimated from historical drop rates;
for clinical trials it can be estimated from prior reporting compliance.
We expect future versions of the crate to provide helpers for these
common estimators, but the abstraction itself is intentionally agnostic.

\subsection{Adjacent constructions this work does not displace}
\label{sec:disc:notdisplace}

Four communities have already addressed pieces of the problem
$\Mtype\langle T \rangle$ solves, and a careful reading should clarify
where this work sits relative to each.

Probabilistic databases~\cite{antova2007maybms,cheng2008orion} factor
tuple existence from attribute value in the relational setting. The
mathematical content is the same as $\Mtype\langle T \rangle$; the
delivery is different. A probabilistic database manages this factoring
inside a query engine, exposing it through SQL-flavoured query languages
and ad-hoc lineage tracking. $\Mtype\langle T \rangle$ exposes the same
factoring as a value-level type usable anywhere
\texttt{Option<T>} would go. Applications that already use a relational
backend should keep using it; applications that compute over values in
the host language now have a type that matches the semantics they need.

Sensor-fusion control theory~\cite{sinopoli2004kalman} factors presence
from value inside specific numerical filters. Sinopoli et al.'s analysis
of Kalman filtering with intermittent observations is mathematically
identical to the Bernoulli-times-Gaussian construction
$\Mtype\langle f64 \rangle$ implements, and the convergence theory there
gives strong guidance for designing filters that survive dropouts.
$\Mtype\langle T \rangle$ does not replace this work; it complements it
by making the factoring available to application code that lives
outside the filter.

Probabilistic programming languages handle missing data at the modelling
layer~\cite{carpenter2017stan,bingham2019pyro,ge2025turing}. Pyro's
\texttt{MaskedDistribution} is the closest mechanical analogue. PPL
encodings let the modeller state the missingness model and run inference
on it; they require the user to commit to a modelling layer and an
inference engine. $\Mtype\langle T \rangle$ does not run inference. It is
a value-level type that supports sampling and SPRT-gated lifting,
suitable for application code where a full PPL would be the wrong tool.

Probabilistic refinement type systems and relational program
logics~\cite{vasilenko2022safecouplings,avanzini2025erhl} verify
quantitative relational properties of probabilistic programs.
$\Mtype\langle T \rangle$ does not aspire to this layer of the stack. It
is a runtime abstraction with measured overhead, not a verification
formalism. A natural follow-up is to specify and verify properties of
programs that use $\Mtype\langle T \rangle$ using one of these logics;
we leave that to future work.

\newpage